% !TeX root = ../../../main.tex
\subsubsection{Evaluation Methodology}
\ChSixHeadingOnly
\paragraph{Quantitative evaluation of extracted timelines}
\label{ch06:sec:evaluation_methodology_metrics}
We evaluate extracted textual time series using four metrics:\\
\textbf{Event match rate}: Proportion of reference events matched via recursive best matching with cosine distance threshold 0.1 using PubMedBert embeddings (Appendix~\ref{ch06:apd:pseudo}).\\
\textbf{Concordance (c-index)}: Probability of correctly ordering random event pairs.\\
\textbf{Median absolute error (MAE)}: Median time discrepancy (in hours) among matched events.\\
\textbf{Area under log-time curve (AULTC)}: Integrated log-transformed discrepancies capturing accuracy across time scales (detailed derivation in Appendix \ref{ch06:apd:ltcdf}).

A clinician qualitative review of annotation quality on \texttt{sepsis-10} is described in Appendix~\ref{ch06:apd:clinician_quality_review}.

% \subsubsection{Clinician evaluation of annotation quality}
% \label{sec:clinician_quality_methods}

% In addition to quantitative metrics, a physician reviewed extracted timelines for the \texttt{sepsis-10} subset.
% For each case report, the clinician assigned:
% (i) a rank ordering (1 = best, 7 = worst) of the annotations (1 manual and 6 LLM-generated), and
% (ii) a categorical assessment (Excellent, Good, Acceptable, or Poor) based on extraction completeness, temporal ordering accuracy, and timing plausibility.

% For annotation quality review, a physician ranked annotations (1=best, 7=worst) and categorized them as Excellent/Good/Acceptable/Poor based on extraction completeness, ordering accuracy, and timing plausibility.


% !TeX root = ../../../main.tex
\begin{table}[htbp]
\caption{Comparison of event match rate (Event), median concordance ($c$), median absolute error (MAE), and area under the log-time curve (AULTC) across LLMs for i2m4, \texttt{sepsis-40}, and \texttt{sepsis-100} datasets.}\ChSixFloatLabel{ch06:tab:results}
\ChSixTableBodyStart
\begin{ChSixFloatBody}

\small


\footnotesize
\centering

\begin{adjustbox}{max width=\linewidth}\begin{tabular}{l*{12}{c}}
\toprule

& \multicolumn{4}{c}{i2m4}
& \multicolumn{4}{c}{\texttt{sepsis-40}}
& \multicolumn{4}{c}{\texttt{sepsis-100}} \\

\cmidrule(lr){2-5}
\cmidrule(lr){6-9}
\cmidrule(lr){10-13}

\textbf{Model}
& \scriptsize Event
& \scriptsize c
& \scriptsize MAE
& \scriptsize AULTC
& \scriptsize Event
& \scriptsize c
& \scriptsize MAE
& \scriptsize AULTC
& \scriptsize Event
& \scriptsize c
& \scriptsize MAE
& \scriptsize AULTC \\

\midrule

GPT-5
& -- & -- & -- & --
& {\ChSixBfTab 0.93} & {\ChSixBfTab 0.965} & {\ChSixBfTab 4} & {\ChSixBfTab 0.822}
& -- & -- & -- & -- \\

DeepSeek-R1
& 0.31 & 0.755 & 13 & 0.786
& 0.79 & 0.892 & 16 & 0.781
& {\ChSixBfTab 0.77} & 0.879 & {\ChSixBfTab 2} & {\ChSixBfTab 0.812} \\

L3.3 70B
& {\ChSixBfTab 0.43} & {\ChSixBfTab 0.768} & {\ChSixBfTab 7} & {\ChSixBfTab 0.805}
& 0.76 & 0.908 & 24 & 0.742
& 0.69 & {\ChSixBfTab 0.920} & 7 & 0.796 \\

L3.1 70B
& 0.32 & 0.747 & 8 & 0.804
& 0.66 & 0.931 & 30 & 0.737
& 0.67 & 0.914 & 6 & 0.794 \\

Mixtral 8x7B
& 0.31 & 0.745 & 21 & 0.773
& 0.68 & 0.906 & 48 & 0.710
& 0.63 & 0.875 & 24 & 0.764 \\

L3.1 8B
& 0.35 & 0.727 & 26 & 0.738
& 0.53 & 0.856 & 84 & 0.685
& 0.66 & 0.772 & 50 & 0.716 \\

BERT 5-CV
& -- & 0.704 & 34 & 0.520
& -- & -- & -- & --
& -- & -- & -- & -- \\

\bottomrule
\end{tabular}\end{adjustbox}

\end{ChSixFloatBody}
\end{table}


